% ID: FIS-TER-EQ-001
% Titolo: Miscela di acqua calda e fredda
% Disciplina: Fisica
% Area: Termologia
% Argomento: Temperatura di equilibrio
% Sottoargomento: Calorimetria senza dispersioni
% Classe: Seconda liceo scientifico
% Difficolta: 3
% Tipo: problema_numerico
% Risultato: 44 ^\circ C
% Tempo_stimato: 8 min
% Competenze: modellizzazione, proporzionalita, conservazione dell'energia
% Prerequisiti: calore specifico, equilibrio termico
% Tag: termologia, calorimetria, temperatura_equilibrio, acqua, conservazione_energia
% Fonte: esempio_originale
% Autore: Riccardo Carli
% Licenza: CC-BY-SA-4.0

\begin{esercizio}
Si mescolano $200\,\mathrm{g}$ di acqua a $20\,^\circ\mathrm{C}$ con
$300\,\mathrm{g}$ di acqua a $60\,^\circ\mathrm{C}$. Trascurando le dispersioni
di calore e la capacita termica del recipiente, determina la temperatura finale
di equilibrio.
\end{esercizio}

\begin{soluzione}
Poiche le due masse sono entrambe acqua, il calore specifico e lo stesso e si
semplifica nel bilancio termico.

Indicando con $T_f$ la temperatura finale:
\[
m_1 c (T_f-T_1)+m_2 c (T_f-T_2)=0.
\]

Semplificando $c$:
\[
200(T_f-20)+300(T_f-60)=0.
\]

\[
200T_f-4000+300T_f-18000=0
\]

\[
500T_f=22000
\qquad\Rightarrow\qquad
T_f=44\,^\circ\mathrm{C}.
\]

La temperatura finale e quindi $44\,^\circ\mathrm{C}$.
\end{soluzione}
